\documentclass[10pt,a4paper]{amsart}
\usepackage{amsmath,amssymb,amsthm,graphicx}
\usepackage{verbatim}
\newenvironment{reference}{\comment}{\endcomment}
%\newenvironment{reference}{}{}
\newenvironment{slogan}{\comment}{\endcomment}
\newenvironment{history}{\comment}{\endcomment}

% For commutative diagrams we use Xy-pic
\usepackage[all]{xy}

% We use 2cell for 2-commutative diagrams.
\xyoption{2cell}
% Theorem environments.
%
\theoremstyle{plain}
\newtheorem{theorem}[subsection]{Theorem}
\newtheorem{proposition}[subsection]{Proposition}
\newtheorem{lemma}[subsection]{Lemma}

\theoremstyle{definition}
\newtheorem{definition}[subsection]{Definition}
\newtheorem{example}[subsection]{Example}
\newtheorem{exercise}[subsection]{Exercise}
\newtheorem{situation}[subsection]{Situation}

\theoremstyle{remark}
\newtheorem{remark}[subsection]{Remark}
\newtheorem{remarks}[subsection]{Remarks}

\numberwithin{equation}{subsection}

% Macros
%
\def\lim{\mathop{\mathrm{lim}}\nolimits}
\def\colim{\mathop{\mathrm{colim}}\nolimits}
\def\Spec{\mathop{\mathrm{Spec}}}
\def\Hom{\mathop{\mathrm{Hom}}\nolimits}
\def\Ext{\mathop{\mathrm{Ext}}\nolimits}
\def\SheafHom{\mathop{\mathcal{H}\!\mathit{om}}\nolimits}
\def\SheafExt{\mathop{\mathcal{E}\!\mathit{xt}}\nolimits}
\def\Sch{\mathit{Sch}}
\def\Mor{\mathop{\mathrm{Mor}}\nolimits}
\def\Ob{\mathop{\mathrm{Ob}}\nolimits}
\def\Sh{\mathop{\mathit{Sh}}\nolimits}
\def\NL{\mathop{N\!L}\nolimits}
\def\CH{\mathop{\mathrm{CH}}\nolimits}
\def\proetale{{pro\text{-}\acute{e}tale}}
\def\etale{{\acute{e}tale}}
\def\QCoh{\mathit{QCoh}}
\def\Ker{\mathop{\mathrm{Ker}}}
\def\Im{\mathop{\mathrm{Im}}}
\def\Coker{\mathop{\mathrm{Coker}}}
\def\Coim{\mathop{\mathrm{Coim}}}

% Boxtimes
%
\DeclareMathSymbol{\boxtimes}{\mathbin}{AMSa}{"02}

%
% Macros for moduli stacks/spaces
%
\def\QCohstack{\mathcal{QC}\!\mathit{oh}}
\def\Cohstack{\mathcal{C}\!\mathit{oh}}
\def\Spacesstack{\mathcal{S}\!\mathit{paces}}
\def\Quotfunctor{\mathrm{Quot}}
\def\Hilbfunctor{\mathrm{Hilb}}
\def\Curvesstack{\mathcal{C}\!\mathit{urves}}
\def\Polarizedstack{\mathcal{P}\!\mathit{olarized}}
\def\Complexesstack{\mathcal{C}\!\mathit{omplexes}}
% \Pic is the operator that assigns to X its picard group, usage \Pic(X)
% \Picardstack_{X/B} denotes the Picard stack of X over B
% \Picardfunctor_{X/B} denotes the Picard functor of X over B
\def\Pic{\mathop{\mathrm{Pic}}\nolimits}
\def\Picardstack{\mathcal{P}\!\mathit{ic}}
\def\Picardfunctor{\mathrm{Pic}}
\def\Deformationcategory{\mathcal{D}\!\mathit{ef}}
\usepackage{cleveref}
\crefname{originalreference}{original reference}{original references}
\Crefname{originalreference}{Original reference}{Original references}
\makeatletter
\newlabel{lemma-finite-Galois}{{09I1}{source}{}{sourcepointer-1}{}}
\newlabel{lemma-finite-Galois@cref}{{[originalreference][0][]09I1}{[1][0][]source}{}{}{}}
\makeatother
\crefname{originalreference}{original reference}{original references}
\Crefname{originalreference}{Original reference}{Original references}
\makeatletter
\newlabel{lemma-normal-and-automorphisms}{{09HS}{source}{}{sourcepointer-2}{}}
\newlabel{lemma-normal-and-automorphisms@cref}{{[originalreference][0][]09HS}{[1][0][]source}{}{}{}}
\makeatother
\crefname{originalreference}{original reference}{original references}
\Crefname{originalreference}{Original reference}{Original references}
\makeatletter
\newlabel{lemma-primitive-element}{{030N}{source}{}{sourcepointer-3}{}}
\newlabel{lemma-primitive-element@cref}{{[originalreference][0][]030N}{[1][0][]source}{}{}{}}
\makeatother
\crefname{originalreference}{original reference}{original references}
\Crefname{originalreference}{Original reference}{Original references}
\makeatletter
\newlabel{lemma-recognize-separable}{{09H3}{source}{}{sourcepointer-4}{}}
\newlabel{lemma-recognize-separable@cref}{{[originalreference][0][]09H3}{[1][0][]source}{}{}{}}
\makeatother
\title{A finite group acting faithfully on a field gives a Galois extension of degree the order of the group over the fixed field}
\author{The Stacks Project Authors}
\date{stacks-project:fields.tex:local-snapshot-20261012; GNU FDL 1.2 or later}
\begin{document}
\maketitle
\noindent\textit{Editorial note.} All delivered mathematical text is the author's own, copied from the pinned Stacks Project chapter file fields.tex; no statement or proof step was rewritten, reordered or removed.\par\medskip
\noindent\textit{Editorial note.} The counted proof is the author's complete argument: The author flags the G-invariance of the orbit polynomial as 'the key to the whole lemma' (source line 2956) and carries the permutation computation in full (source lines 2958-2962).\par\medskip
\noindent\textit{Editorial note.} The label lemma-galois-over-fixed-field names the counted statement, and the proof environment immediately following it in the source is delivered verbatim as its proof.\par\medskip
\noindent\textit{Editorial note.} The counted statement is carried by the author's own lemma environment in the chapter source.\par\medskip
\noindent\textit{Allowed standard prerequisites (editorial).}\begin{itemize}
\item Context span 2935-2941: definition of the fixed field K\textasciicircum{}G
\item Same chapter: lemma-finite-Galois (finite extension Galois iff |Aut(E/F)| = [E:F], stated and proved at source lines 2880-2894), lemma-recognize-separable (a divisor of a separable polynomial is separable), lemma-primitive-element (finite separable extensions are simple), lemma-normal-and-automorphisms (the automorphism bound inequality)
\end{itemize}\medskip
\ifdefined\hypertarget\hypertarget{sourcepointer-1}{}\fi
\noindent\textit{Source reference.} Original reference 09I1: sources/stacks-fields/pinned-original-source.tex, source lines 2881--2881. Prerequisite pointer invoked at 2973; it names a standard fact of the same chapter that the counted argument uses, and every step of the counted proof built on it is retained verbatim in the delivered body.\par\medskip
\ifdefined\hypertarget\hypertarget{sourcepointer-2}{}\fi
\noindent\textit{Source reference.} Original reference 09HS: sources/stacks-fields/pinned-original-source.tex, source lines 2211--2211. Prerequisite pointer invoked at 2997; it names a standard fact of the same chapter that the counted argument uses, and every step of the counted proof built on it is retained verbatim in the delivered body.\par\medskip
\ifdefined\hypertarget\hypertarget{sourcepointer-3}{}\fi
\noindent\textit{Source reference.} Original reference 030N: sources/stacks-fields/pinned-original-source.tex, source lines 2518--2518. Prerequisite pointer invoked at 2991; it names a standard fact of the same chapter that the counted argument uses, and every step of the counted proof built on it is retained verbatim in the delivered body.\par\medskip
\ifdefined\hypertarget\hypertarget{sourcepointer-4}{}\fi
\noindent\textit{Source reference.} Original reference 09H3: sources/stacks-fields/pinned-original-source.tex, source lines 1424--1424. Prerequisite pointer invoked at 2969; it names a standard fact of the same chapter that the counted argument uses, and every step of the counted proof built on it is retained verbatim in the delivered body.\par\medskip
\noindent
Let $G$ be a group acting on a field $K$ (by field automorphisms).
We will often use the notation
$$
K^G = \{x \in K \mid \sigma(x) = x \ \forall \sigma \in G\}
$$
and we will call this the {\it fixed field} for the action of $G$ on $K$.

\begin{lemma}
\label{lemma-galois-over-fixed-field}
Let $K$ be a field. Let $G$ be a finite group acting faithfully on $K$.
Then the extension $K/K^G$ is Galois, we have $[K : K^G] = |G|$,
and the Galois group of the extension is $G$.
\end{lemma}

\begin{proof}
Given $\alpha \in K$ consider the orbit $G \cdot \alpha \subset K$
of $\alpha$ under the group action. Consider the polynomial
$$
P = \prod\nolimits_{\beta \in G \cdot \alpha} (x - \beta) \in K[x]
$$
The key to the whole lemma is that this polynomial is invariant
under the action of $G$ and hence has coefficients in $K^G$.
Namely, for $\tau \in G$ we have
$$
P^\tau = \prod\nolimits_{\beta \in G \cdot \alpha} (x - \tau(\beta)) =
\prod\nolimits_{\beta \in G \cdot \alpha} (x - \beta) = P
$$
because the map $\beta \mapsto \tau(\beta)$ is a permutation of
the orbit $G \cdot \alpha$. Thus $P \in K^G[x]$. Since also
$P(\alpha) = 0$ as $\alpha$ is an element of its orbit
we conclude that the extension $K/K^G$ is algebraic. Moreover,
the minimal polynomial $Q$ of $\alpha$ over $K^G$ divides the
polynomial $P$ just constructed. Hence $Q$ is separable
(by Lemma \ref{lemma-recognize-separable} for example) and
we conclude that $K/K^G$ is separable. Thus $K/K^G$ is Galois.
To finish the proof it suffices to show that $[K : K^G] = |G|$
since then $G$ will be the Galois group by
Lemma \ref{lemma-finite-Galois}.

\medskip\noindent
Pick finitely many elements $\alpha_i \in K$, $i = 1, \ldots, n$
such that $\sigma(\alpha_i) = \alpha_i$ for $i = 1, \ldots, n$ implies
$\sigma$ is the neutral element of $G$. Set
$$
L = K^G(\{\sigma(\alpha_i); 1 \leq i \leq n, \sigma \in G\}) \subset K
$$
and observe that the action of $G$ on $K$ induces an action of $G$ on $L$.
We will show that $L$ has degree $|G|$ over $K^G$. This will finish the
proof, since if $L \subset K$ is proper, then we can add an element
$\alpha \in K$, $\alpha \not \in L$ to our list of elements
$\alpha_1, \ldots, \alpha_n$ without increasing $L$ which is absurd.
This reduces us to the case that $K/K^G$ is finite which is
treated in the next paragraph.

\medskip\noindent
Assume $K/K^G$ is finite. By Lemma \ref{lemma-primitive-element}
we can find $\alpha \in K$ such that $K = K^G(\alpha)$.
By the construction in the first paragraph of this proof we see
that $\alpha$ has degree at most $|G|$ over $K$. However, the
degree cannot be less than $|G|$ as $G$ acts faithfully on
$K^G(\alpha) = L$ by construction and the inequality of
Lemma \ref{lemma-normal-and-automorphisms}.
\end{proof}

\end{document}
